\begin{ExerciseBlock}
\ExerciseHeading{05}{Hund und das Leckerli}{angular_momentum_dog_and_treat}
Auf einem Tisch bewegt sich ein Leckerli (Masse $m$) reibungsfrei mit konstantem Geschwindigkeitsbetrag $v_0$ auf einer Kreisbahn mit Radius $r_0$. Das Leckerli ist an einer Schnur angebracht, das durch ein Loch auf der Tischplatte führt. Ein Hund möchte das Leckerli bekommen, und zieht daher am Faden, um das Leckerli in Richtung Loch zu führen.

\begin{center}
	\CachedTikzFigure{content/resources/chapters/angular_momentum/tikz/dog_and_treat.tex}
\end{center}

\begin{enumerate}
	\item Zeigen Sie: Die Zugkraft des Hundes muss
\[ F(r) = \frac{m v_0^2 r_0^2}{r^3} \]
betragen, um das Leckerli auf einer Kreisbahn mit Radius $r$ (mit $r < r_0$) zu halten.

	\item Würde nach dieser Formel der Hund es jemals schaffen, dass das Leckerli das Loch (also
$r = 0$) erreicht?

	\item Bestimmen Sie die zu verrichtende Arbeit $\Delta W$, um den Kreisbahnradius des Leckerlis auf $r=r_0/2$ zu reduzieren. Drücken Sie die Antwort jeweils durch die kinetische Energie des Leckerlis am Anfang (Kreisbahnradius $r_0$) aus! \\
	Führen Sie die Rechnung auf 2 Arten durch:
	\begin{itemize}
		\item Berechnen Sie die Änderung der Rotationsenergie bzw. des effektiven Potentials.
		\item Berechnen Sie das Arbeitsintegral \[ \Delta W = \int_{A}^{B} \vec{F}(\vec{r}) \cdot d\vec{r} \] für die Zugkraft $F(r)$.
	\end{itemize}

\end{enumerate}
\end{ExerciseBlock}
\UseExerciseSolution{angular_momentum_dog_and_treat}
